% Options for packages loaded elsewhere
\PassOptionsToPackage{unicode}{hyperref}
\PassOptionsToPackage{hyphens}{url}
\PassOptionsToPackage{dvipsnames,svgnames,x11names}{xcolor}
%
\documentclass[
  11pt,
  a4paper,
]{ltjsarticle}
\usepackage{amsmath,amssymb}
\usepackage{iftex}
\ifPDFTeX
  \usepackage[T1]{fontenc}
  \usepackage[utf8]{inputenc}
  \usepackage{textcomp} % provide euro and other symbols
\else % if luatex or xetex
  \usepackage{unicode-math} % this also loads fontspec
  \defaultfontfeatures{Scale=MatchLowercase}
  \defaultfontfeatures[\rmfamily]{Ligatures=TeX,Scale=1}
\fi
\usepackage{lmodern}
\ifPDFTeX\else
  % xetex/luatex font selection
\fi
% Use upquote if available, for straight quotes in verbatim environments
\IfFileExists{upquote.sty}{\usepackage{upquote}}{}
\IfFileExists{microtype.sty}{% use microtype if available
  \usepackage[]{microtype}
  \UseMicrotypeSet[protrusion]{basicmath} % disable protrusion for tt fonts
}{}
\makeatletter
\@ifundefined{KOMAClassName}{% if non-KOMA class
  \IfFileExists{parskip.sty}{%
    \usepackage{parskip}
  }{% else
    \setlength{\parindent}{0pt}
    \setlength{\parskip}{6pt plus 2pt minus 1pt}}
}{% if KOMA class
  \KOMAoptions{parskip=half}}
\makeatother
\usepackage{xcolor}
\usepackage[margin=24mm]{geometry}
\setlength{\emergencystretch}{3em} % prevent overfull lines
\providecommand{\tightlist}{%
  \setlength{\itemsep}{0pt}\setlength{\parskip}{0pt}}
\setcounter{secnumdepth}{5}
\ifLuaTeX
  \usepackage{selnolig}  % disable illegal ligatures
\fi
\IfFileExists{bookmark.sty}{\usepackage{bookmark}}{\usepackage{hyperref}}
\IfFileExists{xurl.sty}{\usepackage{xurl}}{} % add URL line breaks if available
\urlstyle{same}
\hypersetup{
  colorlinks=true,
  linkcolor={blue},
  filecolor={Maroon},
  citecolor={Blue},
  urlcolor={blue},
  pdfcreator={LaTeX via pandoc}}

\author{}
\date{}

\begin{document}

\hypertarget{ux6f14ux7fd2-01-ux30d9ux30afux30c8ux30ebux3068ux30d9ux30afux30c8ux30ebux7a7aux9593}{%
\section{演習 01 ---
ベクトルとベクトル空間}\label{ux6f14ux7fd2-01-ux30d9ux30afux30c8ux30ebux3068ux30d9ux30afux30c8ux30ebux7a7aux9593}}

\hypertarget{ux4e00ux6b21ux7d50ux5408}{%
\subsection{1. 一次結合}\label{ux4e00ux6b21ux7d50ux5408}}

\[
v_1=(1,1,0)^T,\quad v_2=(1,0,1)^T
\]

について、\texttt{b=(3,1,2)\^{}T} が \texttt{span\{v\_1,v\_2\}}
に属するか調べてください。

\hypertarget{ux4e00ux6b21ux72ecux7acb}{%
\subsection{2. 一次独立}\label{ux4e00ux6b21ux72ecux7acb}}

\[
v_1=(1,2,3)^T,\ v_2=(2,4,6)^T,\ v_3=(0,1,1)^T
\]

は一次独立ですか。定義から説明してください。

\hypertarget{ux57faux5e95}{%
\subsection{3. 基底}\label{ux57faux5e95}}

\texttt{R\^{}3} の3本のベクトルが一次独立なら、それらが自動的に
\texttt{R\^{}3} の基底になる理由を説明してください。

\hypertarget{ux90e8ux5206ux7a7aux9593}{%
\subsection{4. 部分空間}\label{ux90e8ux5206ux7a7aux9593}}

\[
W=\{(x,y,z)^T\in\mathbb R^3:x+y+z=0\}
\]

が部分空間であることを示し、基底と次元を求めてください。

\hypertarget{ux6df1ux6398ux308a}{%
\subsection{5. 深掘り}\label{ux6df1ux6398ux308a}}

「座標」と「ベクトルそのもの」はなぜ区別すべきでしょうか。基底を変えたとき何が変化し、何が変化しないか説明してください。

\end{document}
